\chapter{分治}

\section{偏序数点}

在前一章中我们提到过，可以用分治来查询偏序范围最值乃至任意分治信息。我们在此给出做法。

不妨假设范围是 $d$ 个正交半平面围出来的，也就是每个点是一个 $d$ 维带权点 $\mathbf x_i$，每次查询 $\mathbf x_i\le \mathbf y$ 的所有点的分治信息和\graffiti{由于这里没有顺序，需要分治信息可交换}。

我们对第一维进行分治：\texttt{solve(l,r)} 处理所有 $x_1,y_1\in[l,r]$ 的那些点与查询之间的贡献。那么 \texttt{solve(l,r)} 可以拆成三部分：\texttt{solve(l,m)}，\texttt{solve(m+1,r)}，以及处理所有 $x_1\in [l,m]$，$y_1\in [m+1,r]$ 的那些点与查询之间的贡献。注意到第三部分变成了一个相同规模但是维数减一的问题，可以递归到维数减一的问题。来算一下复杂度：设 $T(d,n,V)$ 是维数为 $d$，规模为 $n$，值域大小为 $V$ 时 \texttt{solve} 的复杂度，那么 $T(d,n,V)\le T(d,n_l,V/2)+T(d,n_r,V/2)+T(d-1,n,V_0)+O(n)$。

这里的初值是关键的。考虑 $d=1$ 的情况，如果直接按照分治的方法，那么我们需要 $O(n\log V)$ 的时间；而如果事先排好序了，那么就只需要 $O(n)$ 的时间了。所以，取决于点是否已经排好序了，$T(1,n,V)$ 为 $O(n)$ 或者 $O(n\log n)$。$d=2$ 的时候则是标准的二维数点，可以通过上一章中的扫描线方法做到 $O(n\log n)$。不过，我们在这里也介绍一种纯分治的方法。由于当点排好序时，$T(1,n,V)=O(n\log n)$，我们可以让 \texttt{solve(l,r)} 顺便支持将其内的所有点和询问按照第二维排序。这样，在递归处理完两个子问题之后，就可以 $O(n)$ 计算跨过中点的贡献，并且可以通过归并来 $O(n)$ 排序整个区间。这样我们也得到 $T(2,n,V)=T(2,n_l,V/2)+T(2,n_r,V/2)+O(n)=O(n\log V)$。对于更大的 $d$，使用前面的递归式即可算出 $T(d,n,V)=O(n\log^{d-1}V)$。另外，由于 $V$ 总是可以通过离散化变成 $O(n)$，我们可以在 $O(n\log^{d-1}n)$ 的时间内解决这个问题。

\begin{remark}
    当分治结构与询问无关时，将分治过程记录下来即可得到支持在线询问的数据结构，即树套树。所以，这个问题也存在 $O(\log^{d-1}n)$ 的在线做法，只不过这超出了大纲的范围。
\end{remark}

%TODO：添加算法
在实践中，$d$ 几乎总是不超过 $3$，更大的情况通常需要寻找其他方法\graffiti{比如数点的分块 bitset $O(n^2/w)$，一般情况的 KD tree $O(n^{2-1/d})$}。通常来说，分治方法在 $d=2$ 时的子问题的处理是最关键的。

\begin{exercise}[二维偏序]
    计算一个序列的逆序对数量。
\end{exercise}
\begin{exercise}[三维偏序]
    给定序列 $a,b,c$，查询满足 $a_i\le a_j,b_i\le b_j,c_i\le c_j$ 的 $(i,j)$ 的数量。
\end{exercise}
\begin{exercise}[带修二维数点]
    二维平面，支持插入或删除一个带权点 $\mathbf x$，查询所有 $\mathbf x\le \mathbf y$ 的点的分治信息和。
\end{exercise}
\begin{exercise}[P3157]
    每次删除一个数，查询剩下的序列的逆序对数
\end{exercise}
\begin{exercise}[P4169]
    二维平面，支持插入一个点，查询离一个点最近的点到它的距离，距离为曼哈顿距离 $|\Delta x|+|\Delta y|$。
\end{exercise}

\subsection{半在线}

我们的分治过程事实上可以支持一种比较弱的在线操作：假设点和询问初始全部已知，且在第一维上有序，但点权不是在最开始即给定的，而是在回答询问后才给出。比如说，给定 $(1,\mathbf x_1)$ 的点权，在回答询问 $(2,\mathbf y_2)$ 后再给定 $(3,\mathbf x_3)$ 的点权，然后再在回答询问 $(4,\mathbf y_4)$ 后给定 $(5,\mathbf x_5)$ 的点权，以此类推。由于计算跨过中点的贡献时只需要左边的点权全部已知，我们可以先递归左边，然后处理跨过中点的贡献，再递归右边。这在 OI 里通常被称为\textbf{cdq 分治}\graffiti{OI 中最早是陈丹琦推广的。但是没有人正式定义过 cdq 分治到底指的是什么。}。

\begin{exercise}[P2487]
    二维偏序意义下的 LDS，还需要计数并计数每个点在多少种方案中出现。
\end{exercise}

\begin{exercise}[P4093]
\end{exercise}

